\chapter{Odynophagia}

\section*{Definition}
Odynophagia is pain on swallowing that occurs during the pharyngeal or esophageal phase of deglutition, as distinct from dysphagia which refers to difficulty or sensation of food sticking. The pain is often described as sharp, burning, or pressure-like and can be severe enough to limit oral intake.

\section*{Pathophysiology}
Pain arises from inflammation, ulceration, or mechanical distention of the pharyngeal or esophageal mucosa. Nociceptors in the submucosa are activated by inflammatory mediators (prostaglandins, bradykinin, substance P) released in response to infection, chemical injury, or physical trauma. Esophageal mural tension from obstruction or dysmotility can also trigger visceral afferent pain fibers. Acid exposure in gastroesophageal reflux disease directly damages the epithelium and sensitizes nerve endings.

\section*{Etiology}
\begin{itemize}
  \item Infectious causes: Candida esophagitis (immunocompromised, prolonged antibiotic/steroid use), Herpes simplex virus esophagitis, Cytomegalovirus esophagitis (HIV, transplant), Bacterial esophagitis (Mycobacterium, Nocardia), Severe pharyngitis/tonsillitis (Group A Streptococcus, EBV, gonorrhea), HIV seroconversion
  \item Reflux-related: Severe GERD, Erosive esophagitis, Peptic stricture
  \item Pill-induced esophagitis: Doxycycline, Tetracycline, Bisphosphonates, NSAIDs, Potassium chloride, Quinidine, Iron sulfate
  \item Inflammatory/immune: Eosinophilic esophagitis, Crohn disease involving the esophagus, Lichen planus, Pemphigus vulgaris
  \item Neoplastic causes: Esophageal carcinoma, Esophageal lymphoma, Pharyngeal tumors
  \item Traumatic/mechanical: Foreign body impaction, Post-esophageal instrumentation, Caustic ingestion (acid, alkali), Radiation esophagitis, Nasogastric tube trauma
\end{itemize}

\section*{Indications for Evaluation}
Seek care if:
\begin{itemize}
  \item Odynophagia preventing adequate oral intake or causing weight loss
  \item Pain associated with fever, chills, or neck swelling
  \item Immunocompromised patient with any odynophagia (suspect opportunistic infection)
  \item Known caustic ingestion or suspected foreign body
  \item Odynophagia with hematemesis or melena
\end{itemize}

\section*{Related Symptoms}
\begin{itemize}
  \item Dysphagia
  \item Globus sensation
  \item Regurgitation
  \item Heartburn
  \item Chest pain
  \item Hematemesis
  \item Weight loss
  \item Sore throat
  \item Fever
  \item Oral thrush
\end{itemize}
\textbf{Note:} Endoscopy with mucosal biopsy is essential for definitive diagnosis when odynophagia persists despite empirical treatment or occurs in immunocompromised hosts.

\section*{Self-Care / Home Management}
\begin{itemize}
  \item Eat soft, bland foods (yogurt, pudding, mashed potatoes, smoothies) and avoid spicy, acidic, or coarse foods.
  \item Drink cool liquids or suck on ice chips to soothe irritated mucosa.
  \item Take acetaminophen or NSAIDs for pain relief when appropriate.
  \item Elevate the head of the bed to reduce nocturnal acid reflux.
  \item For suspected pill esophagitis, take medications with a full glass of water and remain upright.
\end{itemize}

\section*{Diagnostic Approach}
\begin{itemize}
  \item History focused on duration, severity, associated symptoms, immune status, and medication list.
  \item Physical exam including oropharyngeal inspection for thrush, ulcers, and tonsillar exudate.
  \item Esophagogastroduodenoscopy (EGD) with biopsy for definitive mucosal evaluation.
  \item Barium swallow study when esophageal stricture, web, or extrinsic compression is suspected.
  \item CT neck/chest for suspected perforation, abscess, or malignancy.
\end{itemize}

\section*{Treatment / Management Overview}
\begin{itemize}
  \item Infectious esophagitis: Antifungals (fluconazole for Candida), antivirals (acyclovir for HSV, ganciclovir for CMV).
  \item Reflux esophagitis: Proton pump inhibitors (PPIs) and lifestyle modifications.
  \item Pill esophagitis: Discontinue offending drug, provide supportive care and PPIs.
  \item Eosinophilic esophagitis: Topical swallowed corticosteroids (fluticasone, budesonide), dietary elimination, PPI trial.
  \item Esophageal dilation for strictures, foreign body removal, or malignancy management.
\end{itemize}

\clearpage
